\chapter{MIDI Routing Setup}
\label{ch:midi_routing}

Duq requires MIDI input to trigger its envelopes. Routing methods vary depending on your DAW.

\section{Ableton Live}
1. Create a MIDI track with your trigger pattern.
2. On the MIDI track, set "MIDI To" to the Audio track containing Duq.
3. In the second dropdown under "MIDI To", select "Duq".
4. Set the MIDI track's Monitor to "In" or "Auto".

\section{Logic Pro}
1. Create an "Audio-Controlled-Unit" instance of Duq. This is done by selecting Duq from the "MIDI-controlled Effects" category in the instrument slot of a Software Instrument track.
2. Use the "Sidechain" input menu in the plugin header to select the audio track you want to process.
3. Logic will route MIDI from the Software Instrument track directly to the plugin.

\section{FL Studio}
1. Open Duq in the Mixer.
2. In the Wrapper settings (gear icon), go to the "Vesta" (Plug) tab.
3. Set an "Input port" (e.g., Port 1).
4. Create a "MIDI Out" channel in the Channel Rack and set its port to 1.
5. Patterns played on the MIDI Out channel will now trigger Duq.

\section{Reaper}
1. Create a track for your MIDI triggers and a track for your audio.
2. Click the "Route" button on the MIDI track and add a new send to the audio track.
3. Ensure MIDI is being sent (Audio: None, MIDI: All).
